\documentclass[10pt,a4paper]{amsart}
\usepackage{amsmath,amssymb,amsthm,graphicx}
% Retained minimal preamble for the Stacks Project chapter sites-cohomology.tex (Cohomology on Sites).
% Verbatim excerpt of the chapter preamble: Xy-pic setup (native preamble lines 14-18) and the native
% theorem declarations with the mathematical macros (native preamble lines 154-217). Copied mechanically;
% nothing added, removed or redefined. The wrapper class and the amsmath/amssymb/amsthm/graphicx base
% packages are supplied by the assembler, so the class/externaldocument part of the native preamble is omitted.
% For commutative diagrams we use Xy-pic
\usepackage[all]{xy}

% We use 2cell for 2-commutative diagrams.
\xyoption{2cell}
% Theorem environments.
%
\theoremstyle{plain}
\newtheorem{theorem}[subsection]{Theorem}
\newtheorem{proposition}[subsection]{Proposition}
\newtheorem{lemma}[subsection]{Lemma}

\theoremstyle{definition}
\newtheorem{definition}[subsection]{Definition}
\newtheorem{example}[subsection]{Example}
\newtheorem{exercise}[subsection]{Exercise}
\newtheorem{situation}[subsection]{Situation}

\theoremstyle{remark}
\newtheorem{remark}[subsection]{Remark}
\newtheorem{remarks}[subsection]{Remarks}

\numberwithin{equation}{subsection}

% Macros
%
\def\lim{\mathop{\mathrm{lim}}\nolimits}
\def\colim{\mathop{\mathrm{colim}}\nolimits}
\def\Spec{\mathop{\mathrm{Spec}}}
\def\Hom{\mathop{\mathrm{Hom}}\nolimits}
\def\Ext{\mathop{\mathrm{Ext}}\nolimits}
\def\SheafHom{\mathop{\mathcal{H}\!\mathit{om}}\nolimits}
\def\SheafExt{\mathop{\mathcal{E}\!\mathit{xt}}\nolimits}
\def\Sch{\mathit{Sch}}
\def\Mor{\mathop{\mathrm{Mor}}\nolimits}
\def\Ob{\mathop{\mathrm{Ob}}\nolimits}
\def\Sh{\mathop{\mathit{Sh}}\nolimits}
\def\NL{\mathop{N\!L}\nolimits}
\def\CH{\mathop{\mathrm{CH}}\nolimits}
\def\proetale{{pro\text{-}\acute{e}tale}}
\def\etale{{\acute{e}tale}}
\def\QCoh{\mathit{QCoh}}
\def\Ker{\mathop{\mathrm{Ker}}}
\def\Im{\mathop{\mathrm{Im}}}
\def\Coker{\mathop{\mathrm{Coker}}}
\def\Coim{\mathop{\mathrm{Coim}}}

% Boxtimes
%
\DeclareMathSymbol{\boxtimes}{\mathbin}{AMSa}{"02}

%
% Macros for moduli stacks/spaces
%
\def\QCohstack{\mathcal{QC}\!\mathit{oh}}
\def\Cohstack{\mathcal{C}\!\mathit{oh}}
\def\Spacesstack{\mathcal{S}\!\mathit{paces}}
\def\Quotfunctor{\mathrm{Quot}}
\def\Hilbfunctor{\mathrm{Hilb}}
\def\Curvesstack{\mathcal{C}\!\mathit{urves}}
\def\Polarizedstack{\mathcal{P}\!\mathit{olarized}}
\def\Complexesstack{\mathcal{C}\!\mathit{omplexes}}
% \Pic is the operator that assigns to X its picard group, usage \Pic(X)
% \Picardstack_{X/B} denotes the Picard stack of X over B
% \Picardfunctor_{X/B} denotes the Picard functor of X over B
\def\Pic{\mathop{\mathrm{Pic}}\nolimits}
\def\Picardstack{\mathcal{P}\!\mathit{ic}}
\def\Picardfunctor{\mathrm{Pic}}
\def\Deformationcategory{\mathcal{D}\!\mathit{ef}}
\usepackage{cleveref}
\crefname{originallemma}{original lemma}{original lemmas}
\Crefname{originallemma}{Original lemma}{Original lemmas}
\makeatletter
\newlabel{lemma-K-flat-resolution}{{06YS}{source}{}{sourcepointer-1}{}}
\newlabel{lemma-K-flat-resolution@cref}{{[originallemma][0][]06YS}{[1][0][]source}{}{}{}}
\makeatother
\crefname{originallemma}{original lemma}{original lemmas}
\Crefname{originallemma}{Original lemma}{Original lemmas}
\makeatletter
\newlabel{lemma-derived-base-change}{{06YY}{source}{}{sourcepointer-2}{}}
\newlabel{lemma-derived-base-change@cref}{{[originallemma][0][]06YY}{[1][0][]source}{}{}{}}
\makeatother
\crefname{originallemma}{original lemma}{original lemmas}
\Crefname{originallemma}{Original lemma}{Original lemmas}
\makeatletter
\newlabel{lemma-pullback-K-flat}{{0G7E}{source}{}{sourcepointer-3}{}}
\newlabel{lemma-pullback-K-flat@cref}{{[originallemma][0][]0G7E}{[1][0][]source}{}{}{}}
\makeatother
\crefname{originallemma}{original lemma}{original lemmas}
\Crefname{originallemma}{Original lemma}{Original lemmas}
\makeatletter
\newlabel{lemma-tensor-product-K-flat}{{07A2}{source}{}{sourcepointer-4}{}}
\newlabel{lemma-tensor-product-K-flat@cref}{{[originallemma][0][]07A2}{[1][0][]source}{}{}{}}
\makeatother
\title{Compatibility of the derived pullback with the derived tensor product}
\author{The Stacks Project Authors}
\date{stacks-project:sites-cohomology.tex:local-snapshot-20261011; GNU FDL 1.2 or later}
\begin{document}
\maketitle
\noindent\textit{Editorial note.} All delivered mathematical text is the author's own, copied from the pinned Stacks Project chapter file sites-cohomology.tex; no statement or proof step was rewritten, reordered or removed.\par\medskip
\noindent\textit{Editorial note.} The proof is the author's complete argument: the representatives are chosen explicitly and the comparison of the two paths is carried out with them, the general case following by the defining property of the representatives.\par\medskip
\noindent\textit{Editorial note.} Every labelled reference occurring inside the counted spans is carried in assembly.external\_references as a non-load-bearing navigation pointer (same-chapter ones at their definition site, cross-chapter ones at the invocation inside the counted span), so no load-bearing mathematics is replaced by a stub. No bibliography citation occurs inside the counted spans.\par\medskip
\noindent\textit{Editorial note.} The checked chapter list for same-name collisions is the set of stacks chapters staged in this repository: chow, cohomology, divisors, etale-cohomology, flat, more-algebra, more-morphisms, morphisms, spaces-more-morphisms, varieties together with this card's chapter sites-cohomology.\par\medskip
\noindent\textit{Editorial note.} The preamble is the chapter's own retained minimal preamble (sources/stacks-sites-cohomology/retained-preamble.tex), a verbatim excerpt of the native preamble: the Xy-pic setup and the native theorem declarations with the mathematical macros.\par\medskip
\noindent\textit{Allowed standard prerequisites (editorial).}\begin{itemize}
\item Derived pullback and its identification with the naive pullback on K-flat complexes with flat terms. Same chapter, the construction of Lf\textasciicircum{}* and the result that it agrees with f\textasciicircum{}* on K-flat complexes with flat terms. Supplies the replacements by which the two paths around the diagram are computed.
\item Derived tensor product and the total complex of a double complex. Same chapter, the derived tensor product computed by K-flat representatives and the Tot construction. Supplies the four objects of the displayed diagram and the functoriality used to compare the two paths.
\end{itemize}\medskip
\ifdefined\hypertarget\hypertarget{sourcepointer-1}{}\fi
\noindent\textit{Source reference.} Original lemma 06YS: sources/stacks-sites-cohomology/pinned-original-source.tex, source lines 3207--3207. Navigation pointer invoked at proof line 3791; the counted span contains the pointer only and the cited result is not reproduced inside it, so the pointer is registered as a non-load-bearing navigation stub in the established form for Stacks material. No bibliography citation is involved.\par\medskip
\ifdefined\hypertarget\hypertarget{sourcepointer-2}{}\fi
\noindent\textit{Source reference.} Original lemma 06YY: sources/stacks-sites-cohomology/pinned-original-source.tex, source lines 3585--3585. Navigation pointer invoked at proof line 3880; the counted span contains the pointer only and the cited result is not reproduced inside it, so the pointer is registered as a non-load-bearing navigation stub in the established form for Stacks material. No bibliography citation is involved.\par\medskip
\ifdefined\hypertarget\hypertarget{sourcepointer-3}{}\fi
\noindent\textit{Source reference.} Original lemma 0G7E: sources/stacks-sites-cohomology/pinned-original-source.tex, source lines 3508--3508. Navigation pointer invoked at proof line 3793; the counted span contains the pointer only and the cited result is not reproduced inside it, so the pointer is registered as a non-load-bearing navigation stub in the established form for Stacks material. No bibliography citation is involved.\par\medskip
\ifdefined\hypertarget\hypertarget{sourcepointer-4}{}\fi
\noindent\textit{Source reference.} Original lemma 07A2: sources/stacks-sites-cohomology/pinned-original-source.tex, source lines 3060--3060. Navigation pointer invoked at proof line 3909; the counted span contains the pointer only and the cited result is not reproduced inside it, so the pointer is registered as a non-load-bearing navigation stub in the established form for Stacks material. No bibliography citation is involved.\par\medskip
\begin{lemma}
\label{lemma-tensor-pull-compatibility}
Let $f : (\Sh(\mathcal{C}), \mathcal{O}_\mathcal{C}) \to
(\Sh(\mathcal{D}), \mathcal{O}_\mathcal{D})$ be a morphism of ringed topoi.
Let $\mathcal{K}^\bullet$ and $\mathcal{M}^\bullet$
be complexes of $\mathcal{O}_\mathcal{D}$-modules.
The diagram
$$
\xymatrix{
Lf^*(\mathcal{K}^\bullet
\otimes_{\mathcal{O}_\mathcal{D}}^\mathbf{L}
\mathcal{M}^\bullet) \ar[r] \ar[d] &
Lf^*\text{Tot}(\mathcal{K}^\bullet
\otimes_{\mathcal{O}_\mathcal{D}}
\mathcal{M}^\bullet) \ar[d] \\
Lf^*\mathcal{K}^\bullet \otimes_{\mathcal{O}_\mathcal{C}}^\mathbf{L}
Lf^*\mathcal{M}^\bullet \ar[d] &
f^*\text{Tot}(\mathcal{K}^\bullet
\otimes_{\mathcal{O}_\mathcal{D}}
\mathcal{M}^\bullet) \ar[d] \\
f^*\mathcal{K}^\bullet \otimes_{\mathcal{O}_\mathcal{C}}^\mathbf{L}
f^*\mathcal{M}^\bullet \ar[r] &
\text{Tot}(f^*\mathcal{K}^\bullet \otimes_{\mathcal{O}_\mathcal{C}}
f^*\mathcal{M}^\bullet)
}
$$
commutes.
\end{lemma}

\begin{proof}
We will use the existence of K-flat resolutions with flat terms
(Lemma \ref{lemma-K-flat-resolution}), we will use that derived pullback
is computed by such complexes (Lemma \ref{lemma-derived-base-change}),
and that pullbacks preserve these properties
(Lemma \ref{lemma-pullback-K-flat}). If we choose such
resolutions $\mathcal{P}^\bullet \to \mathcal{K}^\bullet$
and $\mathcal{Q}^\bullet \to \mathcal{M}^\bullet$, then
we see that
$$
\xymatrix{
Lf^*\text{Tot}(\mathcal{P}^\bullet
\otimes_{\mathcal{O}_\mathcal{D}}
\mathcal{Q}^\bullet) \ar[r] \ar[d] &
Lf^*\text{Tot}(\mathcal{K}^\bullet
\otimes_{\mathcal{O}_\mathcal{D}}
\mathcal{M}^\bullet) \ar[d] \\
f^*\text{Tot}(\mathcal{P}^\bullet
\otimes_{\mathcal{O}_\mathcal{D}}
\mathcal{Q}^\bullet) \ar[d] \ar[r] &
f^*\text{Tot}(\mathcal{K}^\bullet
\otimes_{\mathcal{O}_\mathcal{D}}
\mathcal{M}^\bullet) \ar[d] \\
\text{Tot}(f^*\mathcal{P}^\bullet \otimes_{\mathcal{O}_\mathcal{C}}
f^*\mathcal{Q}^\bullet) \ar[r] &
\text{Tot}(f^*\mathcal{K}^\bullet \otimes_{\mathcal{O}_\mathcal{C}}
f^*\mathcal{M}^\bullet)
}
$$
commutes. However, now the left hand side of the diagram
is the left hand side of the diagram by our choice of
$\mathcal{P}^\bullet$ and $\mathcal{Q}^\bullet$ and
Lemma \ref{lemma-tensor-product-K-flat}.
\end{proof}

\end{document}
